\documentclass[10pt]{article}
% Math Packages
\usepackage{amsmath, mathtools}
\usepackage{amssymb}
\usepackage{amsthm}
\usepackage{amsfonts}
\usepackage{bbm}
\usepackage{breqn}
\usepackage[margin=1in]{geometry}
\usepackage{graphicx}
\usepackage{tikz}
\usetikzlibrary{arrows.meta}
\usetikzlibrary{calc}
\usepackage{forest}
\usepackage{tikz-qtree}
\graphicspath{ {./images/} }
\usepackage{hyperref}
\usepackage[capitalize]{cleveref}
\usepackage[shortlabels]{enumitem}
\usetikzlibrary{arrows,matrix,positioning}
\usepackage{multicol}

\usepackage{listings}
\usepackage{xcolor}

\crefname{problem}{Problem}{Problems}
\Crefname{problem}{Problem}{Problems}

\lstset{
  language=R,
  basicstyle=\ttfamily\footnotesize,
  keywordstyle=\color{blue},
  commentstyle=\color{gray},
  stringstyle=\color{red!70!black},
  framesep=2pt,
  framexleftmargin=0pt, 
  frame=single,
  breaklines=false,
  showstringspaces=false
}
% for the pipe symbol
\usepackage[T1]{fontenc}

% Colorful Notes
\usepackage{color} \definecolor{Red}{rgb}{1,0,0} \definecolor{Blue}{rgb}{0,0,1}
\definecolor{Purple}{rgb}{.5,0,.5} \def\red{\color{Red}} \def\blue{\color{Blue}}
\def\gray{\color{gray}} \def\purple{\color{Purple}}
\newcommand{\rnote}[1]{{\red [#1]}} % \rnote{foo} gives '[foo]' in red
\newcommand{\pnote}[1]{{\purple [#1]}} % \pnote{foo} gives '[foo]' in purple
\newcommand{\bnote}[1]{{\blue #1}} % \bnote{foo} gives 'foo' in blue
\newcommand{\gnote}[1]{{\gray #1}} % \gnote{foo} gives 'foo' in gray
\newcommand{\Max}[1]{{\purple [#1]}} % \bnote{foo} then 'foo' is blue

% Claim numbering (the counter restarts after each proof environment)
\newcounter{claimcount}
\setcounter{claimcount}{0}
\newenvironment{claim}{\refstepcounter{claimcount}\par\addvspace{\medskipamount}\noindent\textbf{Claim \arabic{claimcount}:}}{}
\usepackage{etoolbox}
\AtBeginEnvironment{proof}{\setcounter{claimcount}{0}}
\newenvironment{claimproof}{\par\addvspace{\medskipamount}\noindent\textit{Proof of Claim  \arabic{claimcount}.}}{\hfill\ensuremath{\qedsymbol} \tiny{Claim}

  \medskip}
% Add claim support to cleverref
\crefname{claimcount}{Claim}{Claims}

% Math Environments
\newtheorem{theorem}{Theorem}
\newtheorem{assumption}[theorem]{Assumption}
\newtheorem{lemma}[theorem]{Lemma}
\newtheorem{proposition}[theorem]{Proposition}
\newtheorem{corollary}[theorem]{Corollary}
\newtheorem{question}[theorem]{Question}
\theoremstyle{definition}
\newtheorem*{definition}{Definition}
\newtheorem{remark}[theorem]{Remark}
\newtheorem{example}[theorem]{Example}
\newtheorem{notation}[theorem]{Notation}
\newtheorem{problem}[theorem]{Problem}

% Matrices and Column Vectors. 
\usepackage{stackengine}
\setstackgap{L}{1.0\normalbaselineskip}
\usepackage{tabstackengine}
\setstackEOL{;}% row separator
\setstackTAB{,}% column separator
\setstacktabbedgap{1ex}% inter-column gap
\setstackgap{L}{1.5\normalbaselineskip}% inter-row baselineskip
\let\nmatrix\bracketMatrixstack  %Usage: \nmatrix{1,2,3\4,5,6}
\newcommand\cv[1]{\setstackEOL{,}\bracketMatrixstack{#1}} %usage: \cv{1,2,3}

% Custom Math Coqmmands
\newcommand{\vt}{\vskip 5mm} % vertical space
\newcommand{\fl}{\noindent\textbf} % first line
\newcommand{\Fl}{\vt\noindent\textbf} % first line with space above
\newcommand{\norm}[1]{\left\lVert#1\right\rVert} % norm
\newcommand{\pnorm}[1]{\left\lVert#1\right\rVert_p} % p-norm
\newcommand{\qnorm}[1]{\left\lVert#1\right\rVert_q} % q-norm
\newcommand{\1}[1]{\textbf{1}_{\left[#1\right]}} % indicator function 
\def\limn{\lim_{n\to\infty}} % shortcut for lim as n-> infinity
\def\sumn{\sum_{n=1}^{\infty}} % shortcut for sum from n=1 to infinity
\def\sumkn{\sum_{k=1}^{n}} % shortcut for sum from k=1 to n
\def\sumin{\sum_{i=1}^{n}} % shortcut for sum from i=1 to n
\def\SAs{\sigma\text{-algebras}} % shortcut for $\sigma$-algebras
\def\SA{\sigma\text{-algebra}} % shortcut for $\sigma$-algebra
\def\Ft{\mathcal{F}_t} % time-indexed sigma-algebra (t)
\def\Fs{\mathcal{F}_s} % time-indexed sigma-algebra (s)
\def\F{\mathcal{F}} % sigma-algebra
\def\G{\mathcal{G}} % sigma-algebra
\def\R{\mathbb{R}} % Real numbers
\def\N{\mathbb{N}} % Natural numbers
\def\Z{\mathbb{Z}} % Integers
\def\E{\mathbb{E}} % Expectation
\def\P{\mathbb{P}} % Probability
\def\Q{\mathbb{Q}} % Q probability
\def\dist{\text{dist}} %Text 'dist' for things like 'dist(x,y)'
\newcommand{\indep}{\perp \!\!\! \perp}  %independence symbol
\def\Var{\mathrm{Var}} % Variance
\def\tr{\mathrm{tr}} % trace

% Brackets and Parentheses
\def\[{\left [}
    \def\]{\right ]}
% \def\({\left (}
%   \def\){\right )}

\usepackage{color}
\definecolor{Red}{rgb}{1,0,0}
\definecolor{Blue}{rgb}{0,0,1}
\definecolor{Purple}{rgb}{.5,0,.5}
\def\red{\color{Red}}
\def\blue{\color{Blue}}
\def\gray{\color{gray}}
\def\purple{\color{Purple}}
\definecolor{RoyalBlue}{cmyk}{1, 0.50, 0, 0}
\newcommand{\dempfcolor}[1]{{\color{RoyalBlue}#1}} 
\newcommand{\demph}[1]{\textcolor{RoyalBlue}{\textbf{\slshape #1}}} % Slanted

\usepackage{emoji}
% comment exactly one of the following line to show / hide the solutions
% \newcommand{\solution}[1]{{\purple #1}} % uncomment to show the solutions
\newcommand{\solution}[1]{} % uncomment to hide the solutions




\begin{document}
\begin{center}
  \section*{Math 472: Homework 10}
  \textit{Due Wednesday, April 29}
\end{center}


\begin{problem}[Exercise 10.91]
  Let $Y_{1},\ldots,Y_{20}$ be a random sample of size $n=20$ from a normal
  distribution with unknown mean and known variance $\sigma^{2}=5$. We wish to
  test $H_{0}:\mu=7$ versus $H_{a}:\mu>7$.
  \begin{enumerate}[(a)]
    \item Find the uniformly most powerful test with significance level $.05$.
    \item For the test in part (a), find the power at each of the following
    alternative values for $\mu$: $\mu_{a}=7.5, 8.0, 8.5, $ and $9.0$. 
    \item Sketch a graph of the power function.
  \end{enumerate}
\end{problem}


\begin{problem}
  Let $X_1,\ldots,X_n$ be a random sample from the pdf
  \begin{equation*}
    f(x|\theta)=\theta x^{-2},\quad 0<\theta\leq x<\infty.
  \end{equation*}
  \begin{enumerate}[(a)]
    \item What is a sufficient statistic for $\theta$?
    \item Find the MLE of $\theta$.
  \end{enumerate}
\end{problem}


\begin{problem}
  In a given city it is assumed that the number of auto accidents in a given
  year follows a Poisson distribution. In past years, the average number of
  accidents per year was 15 and this year it was 10. Is it justified to claim
  that the accident rate has dropped?
\end{problem}


\begin{problem}[Exercise 10.21]
  \label{problem:exercise-10.21}
  Shear strength measurements derived from unconfined compression tests for two
  types of soils gave the results shown in the following table (measurements
  in tons per square foot, with $s$ the sample standard deviation).
  \begin{center}
    \begin{tabular}{lccc}
      \hline
      \textbf{Soil Type} & \boldmath$n$ & \boldmath$\overline{y}$ & \boldmath$s$ \\
      \hline
      I   & 30 & 1.65 & 0.26 \\
      II  & 35 & 1.43 & 0.22 \\
      \hline
    \end{tabular}
  \end{center}
  Do the soils appear to differ with respect to average shear strength, at the
  1\% significance level? \textit{[Hint: see example 10.7 in the textbook.]}
  
\end{problem}

\begin{problem}
  Refer to \cref{problem:exercise-10.21}. Construct a 99\% confidence interval
  for the difference in mean shear strengths for the two soil types.
  \textit{[Hint: see example 8.8 in the textbook.]}
  \begin{enumerate}[(a)]
    \item Is the value $\mu_{1}-\mu_{2}=0$ inside the interval?
    \item Based on the interval, should the null hypothesis of
    \cref{problem:exercise-10.21} be rejected? Why?
    \item How does the conclusion compare with your conclusion
    \cref{problem:exercise-10.21}?
  \end{enumerate}
\end{problem}

% \begin{problem}[Exercise 10.97]
%   Let $Y_{1},\ldots,Y_{n}$ be independent and identically distributed random
%   variables with probability mass function $p(y\mid \theta)$ given by
%   \begin{equation*}
%     p(1\mid \theta)  = \theta^{2},
%     \quad
%     p(2\mid \theta)  = 2\theta(1-\theta)
%     \quad \text{and} \quad
%     p(3\mid \theta)  = (1-\theta)^{2},
%   \end{equation*}
%   where $\theta\in (0,1)$. Let $N_{i}$ denote the number of observations equal
%   to $i$ for $i=1,2,3$.
%   \begin{enumerate}[(a)]
%     \item Derive the likelihood function $L(\theta)$ as a function of
%     $N_{1},N_{2},$ and $N_{3}$.
%     \item Find the most powerful test for testing $H_{0}:\theta=\theta_{0}$
%     versus $H_{a}:\theta=\theta_{a}$, where $\alpha_{a}>\alpha_{0}$. Show that
%     your test specifies that $H_{0}$ be rejected for certain values of
%     $2N_{1}+N_{2}.$
%     \item How do you determine the value of $k$ so that the test has nominal
%     level $\alpha$? You do not need to do the actual computation. A clear
%     description of how to determine $k$ is adequate.
%     \item Is the test derived in parts (a)-(c) uniformly most powerful for
%     testing $H_{0}:\theta=\theta_{0}$ versus $H_{a}:\theta>\theta_{0}$? Why or
%     why not?
%   \end{enumerate}
% \end{problem}


\end{document}
% Through the fire and the flames we carry on